% !TEX TS-program = xelatex
% !TEX encoding = UTF-8 Unicode
% AUTO-GENERATED by scripts/build-latex-cv.mjs from canonical site data (src/data/*).
% Do not edit manually; run `npm run build:latex` after updating src/data/.
\documentclass[10pt,a4paper]{article}
\usepackage[a4paper,left=19mm,right=19mm,top=16mm,bottom=18mm,footskip=9mm]{geometry}
\usepackage{iftex}
\ifTUTeX
  \usepackage{fontspec}
\else
  \usepackage[T1]{fontenc}
  \usepackage[utf8]{inputenc}
  \usepackage{lmodern}
\fi
\usepackage{xcolor}
\definecolor{cvaccent}{HTML}{0E7490}
\definecolor{cvmuted}{HTML}{475569}
\usepackage{enumitem}
\usepackage{needspace}
\usepackage{fancyhdr}
\usepackage{lastpage}
\usepackage[unicode,colorlinks=true,urlcolor=cvaccent,linkcolor=cvaccent,pdfborder={0 0 0}]{hyperref}
\hypersetup{
  pdftitle={Curriculum Vitae — Donglin (Dennis) Li},
  pdfauthor={Donglin (Dennis) Li},
  pdfcreator={scripts/build-latex-cv.mjs}
}
\urlstyle{same}
\setlength{\parindent}{0pt}
\setlength{\parskip}{0pt}
\setlength{\emergencystretch}{2em}
\setlength{\headheight}{13pt}
\setlength{\headsep}{3mm}
\renewcommand{\baselinestretch}{1.15}
\raggedbottom
\widowpenalty=10000
\clubpenalty=10000
\displaywidowpenalty=10000
\setlist[itemize]{leftmargin=4.5mm,label=\textbullet,itemsep=2pt,topsep=2pt,parsep=0pt,partopsep=0pt}
\pagestyle{fancy}
\fancyhf{}
\renewcommand{\headrulewidth}{0pt}
\renewcommand{\footrulewidth}{0pt}
\fancyfoot[L]{\color{cvmuted}\fontsize{8.3}{10}\selectfont\href{https://donglinli.com/cv/}{donglinli.com/cv}}
\fancyfoot[C]{\color{cvmuted}\fontsize{8.3}{10}\selectfont CV \textbar{} Updated 22 November 2026}
\fancyfoot[R]{\color{cvmuted}\fontsize{8.3}{10}\selectfont\thepage\,/\,\pageref*{LastPage}}
\newcommand{\CVHeader}[3]{%
  {\fontsize{20}{24}\selectfont\bfseries #1\par}%
  \vspace{2pt}{\fontsize{10.5}{13}\selectfont\color{cvmuted}#2\par}%
  \vspace{4pt}{\fontsize{9.2}{11.5}\selectfont #3\par}%
  \vspace{6pt}{\color{cvaccent}\hrule height 0.8pt}\vspace{6pt}}
\newcommand{\CVSection}[2]{%
  \Needspace{6\baselineskip}%
  \vspace{8pt}\phantomsection\pdfbookmark[1]{#2}{sec-#1}%
  {\fontsize{12}{15}\selectfont\bfseries\color{cvaccent} #2\par}%
  \vspace{2pt}{\color{cvaccent}\hrule height 0.35pt}\vspace{5pt}}
\newcommand{\CVEntry}[3]{%
  \Needspace{4\baselineskip}%
  \vspace{2pt}\noindent\begin{tabular*}{\textwidth}{@{\extracolsep{\fill}}l r@{}}%
    \parbox[t]{0.76\textwidth}{\raggedright\textbf{#1}} & {\color{cvmuted}\fontsize{9.4}{11.5}\selectfont #2}%
  \end{tabular*}\par
  \if\relax\detokenize{#3}\relax\else
    \vspace{1pt}{\color{cvmuted}\fontsize{9.6}{12}\selectfont #3\par}%
  \fi}
\newcommand{\CVText}[1]{\vspace{2pt}{\fontsize{9.6}{12.5}\selectfont #1\par}}
\newcommand{\CVMeta}[1]{\vspace{1.5pt}{\color{cvmuted}\fontsize{9.2}{11.5}\selectfont #1\par}}
\newcommand{\CVEntryEnd}{\vspace{4pt}}

\begin{document}
\CVHeader{Donglin (Dennis) Li}{PhD researcher, Electrical and Electronic Engineering, The University of Manchester}{\href{mailto:hello@donglinli.com}{hello@donglinli.com} \textbar{} \href{https://donglinli.com/}{donglinli.com} \textbar{} \href{https://github.com/Daihaolin201}{GitHub} \textbar{} \href{https://www.linkedin.com/in/donglin-li-55424326b}{LinkedIn} \textbar{} \href{https://orcid.org/0009-0007-2130-4862}{ORCID} \textbar{} \href{https://scholar.google.com/citations?user=mBr54wEAAAAJ}{Google Scholar}}

\CVSection{profile}{Profile}
\CVText{\textbf{Bayesian control for cooperative robots.}}
\CVText{PhD researcher in Electrical and Electronic Engineering at The University of Manchester, developing Bayesian dual control for multi-robot source seeking, with hands-on work in embedded and autonomous robots and research software. Teaches laboratory and tutorial classes in control, robotics and digital electronics.}

\CVSection{education}{Education}
\CVEntry{\href{https://www.manchester.ac.uk/study/postgraduate-research/programmes/list/04620/phd-electrical-and-electronic-engineering/}{Doctor of Philosophy (in progress)}}{Sept. 2024 – Present}{The University of Manchester}
\CVEntryEnd
\CVEntry{BEng (Hons.) in Electronic Engineering — First Class Honours}{Sept. 2021 – Jun. 2024}{The University of Manchester}
\CVEntryEnd

\CVSection{research}{Research}
\CVEntry{Bayesian Dual Control for Cooperative Multi-Robot Source Seeking}{Sept. 2024 – Present}{The University of Manchester · PhD Researcher}
\begin{itemize}
  \item Bayesian dual control for robot teams localising an unknown source in a noisy field: each robot’s motion balances reducing uncertainty in the source estimate against converging on it.
  \item Contributions include dual control with consensus mean fusion and a safety projection (CCC 2026), and trace-weighted, event-triggered fusion for communication-efficient source seeking (IAI 2026, Best Paper Award finalist).
\end{itemize}
\CVEntryEnd

\CVSection{publications}{Publications}
\begin{enumerate}[leftmargin=5.5mm,itemsep=4pt,topsep=2pt,parsep=0pt]
  \item \textbf{D. Li}, Z. Li, H. Ma, C. Cao, Z. Ding, \href{https://cms.amss.ac.cn/conferences/ccc26/program/TuSessions.html\#TuA11-4}{Dual Control for Multi-Robot Source Seeking with Consensus Mean Fusion and Safety Projection}, \emph{45th Chinese Control Conference (CCC 2026), Changsha, China.}\\
  {\color{cvmuted}\fontsize{9.4}{11.5}\selectfont First author \textbar{} Accepted.}
  \item \textbf{D. Li}, Z. Li, J. Wang, H. Ma, Z. Ding, Trace-Weighted Event-Triggered Mean Fusion for Efficient Multi-Robot Source Seeking, \emph{8th International Conference on Industrial Artificial Intelligence (IAI 2026), Shenyang, China.}\\
  {\color{cvmuted}\fontsize{9.4}{11.5}\selectfont First author \textbar{} Accepted. Best Paper Award Finalist.}
  \item H. Ma, Z. Li, J. Wang, C. Cao, \textbf{D. Li}, J. Yang, Z. Ding, \href{https://doi.org/10.1016/j.neucom.2026.134973}{Load Balancing for Large Language Model Inference Serving: A Control and Learning Oriented Survey}, \emph{Neurocomputing, vol. 706, art. 134973, 2026.}\\
  {\color{cvmuted}\fontsize{9.4}{11.5}\selectfont Co-author \textbar{} Published online 3 September 2026. \textbar{} \href{https://doi.org/10.1016/j.neucom.2026.134973}{DOI: 10.1016/j.neucom.2026.134973}}
  \item H. Ma, Z. Li, C. Cao, \textbf{D. Li}, Z. Ding, \href{https://ifac-sms.org/conferences/conferences/IFAC26/program/IFAC26_ContentListWeb_4.html\#thc04_03}{Multi-Objective Pipeline Scheduling for LLM Inference in Heterogeneous DePIN}, \emph{23rd IFAC World Congress (IFAC WC 2026), Busan, Korea.}\\
  {\color{cvmuted}\fontsize{9.4}{11.5}\selectfont Co-author \textbar{} Accepted.}
\end{enumerate}

\CVSection{employment}{Experience}
\CVEntry{Graduate Teaching Assistant}{Sept. 2024 – Present}{The University of Manchester}
\begin{itemize}
  \item Teaches laboratory and tutorial classes across six undergraduate and postgraduate courses in control, robotics and digital electronics, coaching students through controller implementation, robot kinematics and circuit design.
  \item Marks coursework and examinations with written feedback and advises on final-year project demonstrations; contributed to summer 2026 marking and course development and is appointed to EEEN11101 for 2026/27.
\end{itemize}
\CVEntryEnd
\CVEntry{Guest Lecturer \& Experiment Designer (Process Control)}{Jun. 2025 – Dec. 2025}{The University of Manchester}
\begin{itemize}
  \item Designed and built an MSc process-control laboratory: a PLC running PID control, linked to MATLAB/Simulink over OPC for real-time monitoring and analysis.
  \item Delivered guest lectures on the underlying control theory and led the laboratory sessions.
\end{itemize}
\CVEntryEnd

\CVSection{projects}{Projects}
\CVEntry{\href{https://hackabot-2026.com/live/project-4/}{Robots pushing a trolley together}}{Mar. 2026}{Hack-a-Bot 2026, UoM RoboSoc \textbar{} Third prize \textbar{} Team leader}
\CVText{Up to three MONA differential-drive robots must push a trolley into a goal zone around obstacles, sensing only through a shared overhead camera and communicating over Wi-Fi. Perception, planning and formation control run centrally on a laptop as one closed loop.}
\CVText{\textbf{My part: }led the team and the whole software stack: ArUco localisation, clearance-aware A* planning, Pure Pursuit and PID tracking, and the robot link}
\CVMeta{\textbf{Stack: }OpenCV · ArUco · A* · Pure Pursuit · PID · TCP}
\CVEntryEnd
\CVEntry{\href{https://uom-robosoc.github.io/tasks/Task\%202.html}{MONA delivery bot}}{Mar. 2025}{Hack-a-Bot 2025, UoM RoboSoc \& Sony \textbar{} First prize}
\CVText{A 24-hour challenge: a MONA robot must deliver autonomously through a tape maze, Start \ensuremath{\rightarrow} Pickup \ensuremath{\rightarrow} Dropoff, in a single timed run, localised only from an overhead camera stream.}
\CVText{\textbf{My part: }the overhead vision, A* planning and PID control, and the Arduino code for the radio link and motor calibration}
\CVMeta{\textbf{Stack: }OpenCV · ArUco · A* · PID · NRF24L01 · Arduino}
\CVEntryEnd
\CVEntry{\href{https://github.com/Daihaolin201/operator-one}{OperatorOne}}{Mar. 2026}{UK AI Agent Hackathon EP4 \ensuremath{\times} OpenClaw, Imperial College London}
\CVText{A multi-agent system that takes a business idea through product, marketing, sales and operations. A CEO orchestrator runs four specialist OpenClaw agents in stage-gated cycles towards a revenue target.}
\CVText{\textbf{My part: }the CEO orchestrator, the dashboard’s run API and timeline, the GLM-5 model connection with its preflight check, the safety switches for live actions, and the submission documents}
\CVMeta{\textbf{Stack: }OpenClaw · Python · GLM-5 · JSON · GitHub Actions}
\CVEntryEnd
\CVEntry{Line tracking and traffic signs}{Apr. 2023}{NVIDIA PuzzleThon, Manchester Robotics and UoM \textbar{} Second prize}
\CVText{An autonomous PuzzleBot (Jetson Nano, ROS) must drive an unseen track with no human input: follow the line, obey red and green lights and act on direction signs.}
\CVText{\textbf{My part: }colour filtering and masking, line following and sign recognition}
\CVMeta{\textbf{Stack: }ROS · Python · OpenCV · Jetson Nano}
\CVEntryEnd
\CVEntry{\href{https://github.com/Daihaolin201/Smart-TrafficCones}{Smart traffic cones}}{Sept. 2023 – May 2024}{Third-year individual project, The University of Manchester}
\CVText{Robotic traffic cones that deploy themselves: four cone vehicles, one leader and three followers, queue at a set spacing and reproduce the leader’s route, directed from a handheld controller over an 802.15.4 radio network.}
\CVText{\textbf{My part: }all of it: the ESP32-S3 control board, the ESP32 and STM32 firmware, the radio protocol, the control loops and the experiments}
\CVMeta{\textbf{Stack: }ESP32-S3 · STM32 · XBee 802.15.4 · PCB · PID}
\CVEntryEnd
\CVEntry{ESP buggy}{Sept. 2022 – May 2023}{Embedded Systems Project, The University of Manchester \textbar{} Third prize}
\CVText{An autonomous line-following buggy on an STM32F401 (Mbed C++) for a course with slopes, a 180\ensuremath{^\circ} turnaround and a controlled stop, built by a team of five from a custom sensor board up.}
\CVText{\textbf{My part: }circuit design and the sensor board (three revisions in Altium), software, and the system summary in the final report}
\CVMeta{\textbf{Stack: }STM32 · Mbed C++ · Altium PCB · PID · BLE}
\CVEntryEnd
\CVEntry{CSF Pilot and research tools}{Jul. 2026 – present}{Desktop software for HPC, literature and figures}
\CVText{CSF Pilot is a desktop client (PySide6) for the University of Manchester CSF3 HPC cluster. It keeps one MFA-authenticated SSH session, multiplexes file transfer, Slurm job control, an editor, a terminal and usage analytics over it, and exposes the session to AI agents through an MCP server.}
\CVMeta{\textbf{Stack: }PySide6 · paramiko · Slurm · MCP · pytest}
\CVEntryEnd
\CVEntry{Teaching workbenches}{Sept. 2026 – present}{Local-first browser apps for tutoring and student support}
\CVText{A local-first, single-file web application that runs a tutoring practice end to end: students, scheduling, lesson preparation, homework, feedback, credit-hour packages and pay reconciliation. The system of record is one JSON file in a user-chosen folder, accessed through the File System Access API.}
\CVMeta{\textbf{Stack: }HTML · JavaScript · File System Access API · ICS · node:test}
\CVEntryEnd

\CVSection{awards}{Awards}
\begin{itemize}[leftmargin=4.5mm,label=\textbullet,itemsep=2.5pt,topsep=2pt,parsep=0pt]
  \item \textbf{2026} Best Paper Award Finalist — IAI 2026 {\color{cvmuted}--- 8th International Conference on Industrial Artificial Intelligence}
  \item \textbf{2026} Third Prize — Hack-a-Bot Robotics Challenge {\color{cvmuted}--- UoM RoboSoc}
  \item \textbf{2026} \href{https://www.se.manchester.ac.uk/doctoral-academy/support/finance-and-funding/doctoral-training-and-development-support-fund/}{Doctoral Training and Development Support Fund (DTDSF)} {\color{cvmuted}--- The University of Manchester}
  \item \textbf{2025} First Place Poster Prize — EEE Postgraduate Research Conference {\color{cvmuted}--- The University of Manchester}
  \item \textbf{2025} First Prize — Hack-a-Bot Robotics Challenge {\color{cvmuted}--- UoM RoboSoc \& Sony}
  \item \textbf{2023} Second Prize — NVIDIA PuzzleThon
  \item \textbf{2023} Third Prize — ESP Buggy Development
\end{itemize}

\CVSection{service}{Service}
\CVEntry{Peer Reviewer}{May–Oct. 2026}{TIMC · IEEE TAES · JACA 2026}
\begin{itemize}
  \item Reviewed five manuscripts: one each for IEEE Transactions on Aerospace and Electronic Systems and TIMC, and three for JACA 2026.
\end{itemize}
\CVEntryEnd
\CVEntry{EEE PGR Representative}{Oct. 2024 – Present}{The University of Manchester}
\begin{itemize}
  \item Leads the department’s PGR representatives for 2025/26, consolidating feedback from doctoral researchers and teaching assistants and raising it with the department and the Doctoral Academy.
\end{itemize}
\CVEntryEnd
\CVEntry{Student \& PGR Ambassador}{Oct. 2024 – Present}{The University of Manchester}
\begin{itemize}
  \item Represents the University at open days and outreach events, presenting engineering programmes and student life and running robotics demonstrations in the laboratories.
\end{itemize}
\CVEntryEnd

\CVSection{skills}{Skills}
\begin{itemize}[leftmargin=4.5mm,label=\textbullet,itemsep=2.5pt,topsep=2pt,parsep=0pt]
  \item \textbf{Research Methods:} Bayesian estimation and dual control; experimental design; Monte Carlo and ablation studies; statistical analysis and optimisation; literature synthesis, academic writing and peer review.
  \item \textbf{Programming:} C/C++; Python (NumPy, SciPy, pandas, Matplotlib); MATLAB; R; JavaScript/TypeScript; SQL; PIC18 assembly; PLC (Ladder Logic, SCL); Bash; HTML/CSS; Linux.
  \item \textbf{Control \& Robotics:} PID, distributed and consensus control; formation control; active source seeking; A* and Pure Pursuit; multi-robot coordination; ROS; OpenCV/ArUco visual localisation.
  \item \textbf{Modelling \& Simulation:} MATLAB/Simulink; LabVIEW/NI myDAQ; PLECS; Multisim; Tanner T-Spice/SPICE; ANSYS.
  \item \textbf{Hardware \& Embedded Systems:} STM32/Mbed; PIC18F8722; ESP32/ESP8266; Arduino; Xilinx FPGA/VHDL; PCB prototyping; bare-metal and RTOS firmware.
  \item \textbf{Teaching \& Assessment:} Laboratory and tutorial teaching; control and robotics demonstrations; coursework and examination marking; technical feedback on student projects.
\end{itemize}

\end{document}
